
Latency es la diferencia entre el \textit{time step} y el periodo, por lo que describe el desfase entre ejecuciones. Aquí se pueden comparar los resultados que muestra la página oficial de \textit{RTXI} con los obtenidos. Esta vez, la que mejores resultados ha obtenido es la prioridad $80$ (Figura \ref{FIG:DISTLAT80}). En este caso, es normal que dé resultados peores, pues el stress test que se ha realizado es bastante más agresivo, usando los parámetros que se usa en su script se obtienen resultados que sí respetan la restricción de los $100$ $\mu s$ (Figura \ref{FIG:DISTLAT90RTXI}), además de tener una media de $0.15$ $\mu s$. 

Cabe destacar que es de esperar que bajo esta prueba se obtengan algunos valores mucho más altos que los que muestra \textit{RTXI} en su página, pues en este caso no se está corriendo un simple comando sino la aplicación entera con \textit{realtime threads}. A pesar de eso, la media y la desviación típica son menores que en las pruebas que ellos muestran. 

\begin{figure}[Distribución de latency con prioridad 80]{FIG:DISTLAT80}{Distribución de \textit{latency} con prioridad $80$.}
    \image{0.7\textwidth}{}{latency-distribution-priority-80}
\end{figure}

\begin{figure}[Distribución de latency con prioridad 90]{FIG:DISTLAT90RTXI}{Distribución de \textit{latency} con prioridad $90$ usando los parámetros de stress del script de \textit{RTXI}.}
    \image{0.7\textwidth}{}{latency-distribution-rtxi-test}
\end{figure}

